\paragraph{Extension to sparse polynomial factors with growing scopes.}
The same integral construction applies beyond constant arity.
Fix integers $C,a\ge1$.  Suppose there are $n\ge2$ Boolean variables
and $m\le Cn$ factors, each given by an explicit monomial list
\[
 \Phi_f(x)=\sum_{j=1}^{\ell_f}y_{f,j}x^{\alpha(f,j)},
 \qquad d_f=\max_j|\alpha(f,j)|,
 \qquad \sum_f d_f\le Cn.
\]
The coefficient symbols $y_{f,j}$ are independent; repeated
monomials in a list are permitted.  An empty list means the zero
factor and is handled immediately.  Let $L$ be the full bit length
of these lists, including their variable indices and binary
exponents.  Suppose each declared scope $I_f$ has size at most
$\lambda^a$, where $\lambda=\lceil\log_2 n\rceil$.
The degree bound is imposed on these explicitly supplied
polynomials before Boolean reduction.  Then
\[
 Z_{\Phi}(y)=\sum_{x\in\{0,1\}^n}\prod_f\Phi_f(x)
\]
has an integral division-free circuit of size and deterministic
bit construction time
\[
 \operatorname{poly}(n,L)\,
 2^{O_{C,a}(n(\log\log n)^2/\sqrt{\log n})}.
\]
Its integer constants have polynomial bit length in $n,L$.
It admits the same constant-free compilation of formal degree
at most $\max(m,1)$, with subtraction gates assigned the
addition degree rule.  The output is the formal polynomial
over $\mathbb Z[y]$, so this single circuit identity transfers
to every commutative coefficient ring.

In particular, fixed-degree factors with a polylogarithmic
number of monomials satisfy these hypotheses after changing
the fixed constants: if the degree is at most $k$ and the
list length at most $(\log n)^b$, the scope is at most
$k(\log n)^b$, and the total degree is at most $kCn$.
This preserves the original number $n$ of Boolean variables.
There are no auxiliary verifier variables, and no truth table
on the entire growing scope is constructed.

We prove the exact quantitative interface.  Append one
dummy factor equal to one, with declared scope $\{i\}$, for
each variable $i$.  These factors contribute no coefficient
symbols and have input degree zero; they only ensure that
every variable has a declared incidence.  If $M_0$ is the
resulting total scope length, then
\[
              n\le M_0\le n+Cn\lambda^a.
\]
Put
\[
 B_0=\max(1,\lceil M_0/n\rceil)\lceil\sqrt\lambda\rceil,
 \qquad H=\{i:\deg_{\rm inc}(i)>B_0\}.
\]
Double-counting incidences gives
$|H|\le M_0/B_0\le n/\sqrt\lambda$.
Sum over all $2^{|H|}$ assignments to these variables.  On a
branch, substitute the chosen bits in the explicit monomial
lists and remove them from each declared scope.  Keep the
remaining declared scopes even if some monomial becomes
zero.  This uses only the syntactic list and chosen bits,
with no test of a formal polynomial identity.

Let $n'=n-|H|$.  For sufficiently large $n$ one has
$n/2\le n'\le n$.  Each surviving variable has incidence at
most $B_0$, every scope still has size at most $\lambda^a$,
and the dummy factors ensure $n'\le M'\le B_0n'$ for the
remaining total incidence $M'$.  Factors with empty scope
can be multiplied outside the branch's counting routine;
their polynomials in the $y$ variables are explicit linear
forms.  This removal does not change $M'$.

Choose a fixed integer $\beta$ larger than $a+2$, increasing
the threshold of the exceptional finite range to absorb
$C$.  With $\lambda'=\lceil\log_2 n'\rceil$, both the maximum
scope and $B_0$ are at most $(\lambda')^\beta$ for all
sufficiently large $n$.  E Lemma~4.8, with
$\nu'=\lceil\log_2\lambda'\rceil$, therefore supplies
\[
 r=\left\lfloor\frac{\sqrt{\lambda'}}{(\nu')^2}\right\rfloor,
 \quad Q=\left\lfloor\frac{\lambda'}{128\nu'}\right\rfloor,
 \quad R=(\beta+3)\lceil\sqrt{\lambda'}\rceil,
 \quad \tau=\frac1{r\lambda'},\quad D=r^{2Q}.
\]
Its exact numerical checks and E Lemma~4.7 construct the
layout in polynomial time.  Input sizes failing those
checks form a finite range depending only on the fixed
constants; direct expansion handles that range within
the asserted asymptotic bound.

The independent-incidence degree in E Section~3.1 is bounded
by the sum of the factor degrees, because replacing the
same Boolean variable in different factors by distinct
indeterminates does not increase any factor's total degree.
Consequently, on every branch,
\[
       d\le\sum_f d_f\le Cn\le 2Cn'.
\]
This is the crucial degree hypothesis; the possibly larger
number $M'$ of incidences is not substituted for $d$.
The interpolation polynomial has degree at most $dR$.
Taking $q=\lfloor dR/Q\rfloor$ and $J=q+1$ gives $QJ>dR$.
E Proposition~3.4 has the explicit enumeration bound
\[
       2Qr\,C_q(2D+1)^{2\tau n'}2^{2n'/r},\qquad
       C_q=\sum_{s=0}^{\min(q,n')}2^s\binom{n'}s.
\]
Put $\eta=R/Q$.  The source's finite-difference estimate
$\eta^qC_q\le(1+2\eta)^{n'}$ yields
\[
 \log_2 C_q
 \le q\log_2(Q/R)+n'\log_2(1+2R/Q)
 =O_C\!\left(n'\frac RQ\log\frac{2Q}{R}\right).
\]
Together with E Lemma~4.8 this bounds all three exponential
costs by
\[
 O_{C,a}\!\left(\frac{n'(\nu')^2}{\sqrt{\lambda'}}\right).
\]
The additional branching cost has logarithm
$|H|\le n/\sqrt\lambda$.  Since $n'=\Theta(n)$, their sum
is $O_{C,a}(n(\log\log n)^2/\sqrt{\log n})$.

It remains to verify the factor and coefficient costs,
which an arity statement alone would not justify.  After
the fixed Boolean substitutions in E Lemma~3.3, evaluate
each listed monomial directly in the truncated algebra
$A[z]/(z^J)$, using its binary exponents or at most
$\sum_i\alpha(f,j)_i\le Cn$ multiplications.  Sum these
values with their formal coefficients $y_{f,j}$.
This takes polynomially many algebra operations in
$n,J,L$; it never enumerates the $2^{|I_f|}$ local
assignments.  Before multiplication by the formal
coefficients, each such substituted monomial is a
scalar Taylor polynomial.  The rows and unit inverses
have the polynomial coordinate heights already proved
for the cyclotomic construction.  A product of at most
$Cn$ such scalar inputs, followed by Taylor truncation,
still has polynomial bit height by the same multiplicative
coordinate-norm bound.  Adding at most $L$ listed terms
increases logarithmic norm by at most $O(\log(L+1))$.
All local scalar preprocessing is therefore polynomial
in $n,L$, and $Q,R,J,D$ are polynomially bounded in $n$.

Use the integral swap matrix, integer measurement fibers,
unit-separated cyclotomic nodes and explicit Hasse
interpolation proved above.  Their correctness depends on
the layout and the degree estimate, not on constant scope
size.  They give formal identities over $A[y]$ at each
stage.  The same coefficient-independent choices of
representatives and branches allow circuit unrolling,
and taking the coordinate of one in $A$ descends the
formal output to $\mathbb Z[y]$.  Each genuine factor
remains linear in its own coefficient block.  Group
products, sums, empty-scope multiplication and the final
sum over branches retain the same block degree.  The
constant-free compilation therefore has formal degree
at most $\max(m,1)$, just as in the constant-arity case.
This finishes the construction and its bit bound.

The hypotheses delimit the extension.  A general truth
table on $\operatorname{polylog} n$ variables can have
interpolation degree equal to its scope size; $O(n)$ such
factors need not have total degree $O(n)$.  In that case
the actual $q=\lfloor dR/Q\rfloor$ must be retained and
the displayed bound has not been proved.  Nor does a
coefficient-dependent compression of a numerical equation
system automatically preserve this universal formal
polynomial in independent parameters.

\paragraph{Geometric statement of the extended trace theorem.}
Fix $C,a\ge1$.  Let $m\le Cn$ explicit polynomial lists
$\Phi_f=\sum_j y_{f,j}x^{\alpha(f,j)}$ have description
length $L$, declared scopes of size at most
$(\lceil\log_2n\rceil)^a$, and total degree budget
$\sum_f\max_j|\alpha(f,j)|\le Cn$.  Put $B=\mathbb Z[y]$ and
\[
 \mathscr A=B[x_1,\ldots,x_n]/(x_i^2-x_i),\qquad
                       u=\prod_f\Phi_f(x).
\]
The finite locally free algebra $\mathscr A$ has rank $2^n$.
Its multiplication trace
$\operatorname{Tr}_{\mathscr A/B}(u)=Z_\Phi(y)$, as is
seen in the orthogonal Boolean-idempotent basis.  This
regular function on $\operatorname{Spec}B$ has an integral
circuit of size and bit construction time
$\operatorname{poly}(n,L)2^{O_{C,a}(n(\log\log n)^2/\sqrt{\log n})}$,
with polynomial-bit integer constants and a constant-free
compilation of formal degree at most $\max(m,1)$, with $n$
the original number of Boolean coordinates.  For every map
$B\to R$ the same circuit computes the base-changed
multiplication trace over $R$; coordinate descent from
the finite free cyclotomic extension proves the identity
before any specialization.  The exact classical trace
and base-change interface is EGA IV4, 18.2.1--18.2.2,
printed p.~111, including equation 18.2.1.2.  This is an
effective theorem for these specified regular trace
functions, not a claimed polynomial-time algorithm for
general finite algebra traces or GCT multiplicities.
